% GENERATED FILE. Edit canonical lessons, then run npm run generate:books.
% canonical-source-hash: fnv1a64:b5ef3f8c11ce5a0d
% canonical-lessons: ML-C104-one-present, ML-C104-ippol, ML-C104-no-am, ML-C104-eppol-answer, ML-R104-present-recall

\chapter{One Present, Two Englishes}
\label{ch:ml-present}
\hlchaptermodality{104}

\begin{chapteropening}
\textbf{By the end of this chapter:} \emph{I can say what I am doing right now, and ask somebody when they are coming.}

From cold: the one present form, the word that pins it to this moment, and the helper Malayalam does not have.
\end{chapteropening}

\section[ñān pōkunnu]{\ml{ഞാൻ} \ml{പോകുന്നു} — I go, and I am going}

\label{lesson:ML-C104-one-present}



\begin{quote}
Say \emph{I go}. Then say it for \emph{you} and for \emph{they} — and notice that the verb did not move.
\end{quote}

\subsection*{What to know first}
You were told that Malayalam's present tense is \textbf{one word for everybody}. Here is the part that was not said:

\textbf{It is also one word for two English tenses.}

\noindent
\begin{tabularx}{\linewidth}{@{}>{\raggedright\arraybackslash}X>{\raggedright\arraybackslash}X>{\raggedright\arraybackslash}X@{}}
\toprule
\textbf{} & \textbf{} & \textbf{} \\
\midrule
\textbf{\ml{ഞാൻ} \ml{പോകുന്നു}} & \emph{ñān pōkunnu} & I go \\
\textbf{\ml{ഞാൻ} \ml{പോകുന്നു}} & \emph{ñān pōkunnu} & I am going \\
\bottomrule
\end{tabularx}

Those are the \textbf{same} Malayalam sentence, written twice. English asks you to choose — \emph{I go to work} against \emph{I am going to work} — and the choice does not exist here. There is nothing to get right.

\begin{culture}
\textbf{English is the one doing something extra here.} It splits its present into a habitual and an ongoing form; Malayalam does not. The ending \textbf{-\ml{ഉന്നു}} covers what you do every day and what you are doing at this moment, and a listener takes whichever the situation offers.

\textbf{This is not a gap in Malayalam.} When it matters which you mean, there is a way to say so — and the one this book gives you is a word in front, not a change to the verb. That is the next lesson.
\end{culture}

\subsection*{Your turn}
Say these aloud:

\begin{itemize}
\raggedright
  \item \emph{ñān pōkunnu} — and mean \emph{I go}
  \item \emph{ñān pōkunnu} again — and mean \emph{I am going}
  \item the same for \emph{to come} — \emph{ñān varunnu}
  \item what English makes you choose that Malayalam does not
\end{itemize}

\begin{tcolorbox}[breakable,colback=teal!4,colframe=teal!35!black,title={Before you move on}]
How do you say \emph{I go}? (\textbf{\ml{ഞാൻ} \ml{പോകുന്നു}}.) How do you say \emph{I am going}? (\textbf{The same} — \ml{ഞാൻ} \ml{പോകുന്നു}.) Which language is doing the extra work here? (\textbf{English} — it splits a present that Malayalam keeps whole.)
\end{tcolorbox}
\section[ippōḷ ñān pōkunnu]{\ml{ഇപ്പോൾ} \ml{ഞാൻ} \ml{പോകുന്നു} — right now, I am going}

\label{lesson:ML-C104-ippol}



\begin{quote}
Say \emph{I am going}. Then say the word for \textbf{now}.
\end{quote}

\subsection*{What to know first}
The last lesson left one thing open: if \textbf{\ml{ഞാൻ} \ml{പോകുന്നു}} can mean either, how do you say you mean \textbf{this moment}?

\textbf{You put the moment in front of it.}

\begin{quote}
\textbf{\ml{ഇപ്പോൾ} \ml{ഞാൻ} \ml{പോകുന്നു}.} — \emph{ippōḷ ñān pōkunnu} — \textquotedblleft{}\textbf{Right now} I am going.\textquotedblright{}
\end{quote}

\textbf{The verb did not change.} Not one letter of it. What changed is that a word naming the moment is standing at the front, and it is a word you have owned for a long while.

\begin{grammarlens}[title={Grammar Lens: the slot again}]
That front position is not new either. \textbf{\ml{ഇന്ന്}} has stood there, and \textbf{\ml{നാളെ}}, and a named day, and an hour with an ending on it. \textbf{\ml{ഇപ്പോൾ}} is the same slot doing the narrowest version of its job: not today, not this Monday — \emph{now}.

\textbf{This is how Malayalam answers a question English answers with the verb.} English changes \emph{go} into \emph{am going}. Malayalam leaves the verb alone and adds a word — and a word is easier than a form.
\end{grammarlens}

\subsection*{Your turn}
Say these aloud:

\begin{itemize}
\raggedright
  \item \emph{ñān pōkunnu}, then \emph{ippōḷ ñān pōkunnu}
  \item the same with \emph{to come} — \emph{ippōḷ ñān varunnu}
  \item which part of the sentence changed, and which did not
\end{itemize}

\begin{tcolorbox}[breakable,colback=teal!4,colframe=teal!35!black,title={Before you move on}]
How do you say \emph{right now I am going}? (\textbf{\ml{ഇപ്പോൾ} \ml{ഞാൻ} \ml{പോകുന്നു}}.) What happened to the verb? (\textbf{Nothing} — it is unchanged.) Where did \textbf{\ml{ഇപ്പോൾ}} go? (\textbf{At the front}, in the slot that says \emph{when}.)
\end{tcolorbox}
\section[ippōḷ ñān varunnu]{\ml{ഇപ്പോൾ} \ml{ഞാൻ} \ml{വരുന്നു} — and there is no word for \emph{am}}

\label{lesson:ML-C104-no-am}



\begin{quote}
Say \emph{to come}. Then say \emph{right now I am going}.
\end{quote}

\subsection*{What to know first}
Take a second verb through it and watch nothing bend:

\begin{quote}
\textbf{\ml{ഇപ്പോൾ} \ml{ഞാൻ} \ml{വരുന്നു}.} — \emph{ippōḷ ñān varunnu} — \textquotedblleft{}Right now I am coming.\textquotedblright{}
\end{quote}

\textbf{You have said that verb before.} \emph{varunnu} was in the tense table you were given for \emph{varuka}, and you drilled it there beside \emph{vannu} and \emph{varuṁ}. What is new is not the word — it is that it does the same double duty \textbf{\ml{പോകുന്നു}} does.

\begin{grammarlens}[title={Grammar Lens: the word you must not reach for}]
English says \emph{I \textbf{am} going}, in two words. When you go looking for the Malayalam word to put where \emph{am} is, you will find \textbf{\ml{ആണ്}} — the copula you were given early on for \emph{is}, and which you have since heard carrying \emph{am} as well.

\textbf{Do not put it here.}

\textbf{\ml{ഇപ്പോൾ} \ml{ഞാൻ} \ml{വരുന്നു}} is right: one verb, nothing added. Putting a copula in the middle of it — \emph{ippōḷ ñān āṇŭ varunnu} — \textbf{is ungrammatical}, and it is written here in sound only, so that the shape never sits on the page looking like Malayalam you could copy.

\textbf{\ml{ആണ്}} joins a thing to what it \emph{is} — a name, a description. That is its whole job, and propping up another verb is not part of it. The whole of \emph{am going} already lives inside \textbf{\ml{വരുന്നു}}; there is no second word for it to be.

\textbf{Malayalam does put verbs together in other places} — you have met \textbf{\ml{ഉണ്ട്}} and \textbf{\ml{കഴിയും}} doing exactly that. This is not one of those places, and \textbf{\ml{ആണ്}} is not one of those verbs.
\end{grammarlens}

\subsection*{Your turn}
Say these aloud:

\begin{itemize}
\raggedright
  \item \emph{ippōḷ ñān varunnu}
  \item \emph{ippōḷ ñān pōkunnu}, then the coming one again
  \item where English puts \emph{am}, and what Malayalam puts there instead
\end{itemize}

\begin{tcolorbox}[breakable,colback=teal!4,colframe=teal!35!black,title={Before you move on}]
\emph{Right now I am coming}? (\textbf{\ml{ഇപ്പോൾ} \ml{ഞാൻ} \ml{വരുന്നു}}.) How many words carry \emph{am coming}? (\textbf{One} — \ml{വരുന്നു}.) Does \textbf{\ml{ആണ്}} belong in that sentence? (\textbf{No} — it joins a thing to what it is; it does not help another verb.)
\end{tcolorbox}
\section[Practice]{\ml{നിങ്ങൾ} \ml{എപ്പോൾ} \ml{പോകും}? — asked long ago, answered now}

\label{lesson:ML-C104-eppol-answer}



\begin{quote}
Say \emph{now}, \emph{then}, and \emph{when?} — the three you were given as a set.
\end{quote}

\subsection*{What to know first}
\textbf{Nothing about this pair is new to you.} You were taught the three beginnings as a system — \textbf{\ml{ഇ}-} near, \textbf{\ml{അ}-} far, \textbf{\ml{എ}-} a question — and then handed this row of it whole:

\noindent
\begin{tabularx}{\linewidth}{@{}>{\raggedright\arraybackslash}X>{\raggedright\arraybackslash}X>{\raggedright\arraybackslash}X@{}}
\toprule
\textbf{} & \textbf{} & \textbf{} \\
\midrule
\textbf{\ml{ഇപ്പോൾ}} & \emph{ippōḷ} & now \\
\textbf{\ml{അപ്പോൾ}} & \emph{appōḷ} & then \\
\textbf{\ml{എപ്പോൾ}} & \emph{eppōḷ} & when? \\
\bottomrule
\end{tabularx}

\textbf{What is new is the answer.} You were given the question \textbf{\ml{നിങ്ങൾ} \ml{എപ്പോൾ} \ml{പോകും}?} with a \textbf{future} on the end of it, and no way to reply except with another future.

\subsection*{Your turn}
Now you can answer it with \textbf{this moment} instead:

\begin{quote}
— \textbf{\ml{നിങ്ങൾ} \ml{എപ്പോൾ} \ml{പോകും}?} — \emph{niṅṅaḷ eppōḷ pōkuṁ?} — \textquotedblleft{}When will you go?\textquotedblright{} — \textbf{\ml{ഇപ്പോൾ} \ml{ഞാൻ} \ml{പോകുന്നു}.} — \emph{ippōḷ ñān pōkunnu.} — \textquotedblleft{}I'm going right now.\textquotedblright{}
\end{quote}

Say these aloud:

\begin{itemize}
\raggedright
  \item the question — \emph{niṅṅaḷ eppōḷ pōkuṁ?}
  \item the answer — \emph{ippōḷ ñān pōkunnu}
  \item the same exchange with \emph{to come} — answer \emph{ippōḷ ñān varunnu}
  \item which word in your answer says \emph{this moment}
\end{itemize}

\begin{grammarlens}[title={What you've built}]
\textbf{The question was asked in the future and answered in the present, and that is not a mismatch} — it is the ordinary way to say you are on your way. The question asks about a plan; the answer says the plan is happening now.

And the answer needs \textbf{one} verb to do it. No word went in where English puts \emph{am}.
\end{grammarlens}

\begin{tcolorbox}[breakable,colback=teal!4,colframe=teal!35!black,title={Before you move on}]
\emph{When?} (\textbf{\ml{എപ്പോൾ}}.) Its two partners? (\textbf{\ml{ഇപ്പോൾ}} \emph{now}, \textbf{\ml{അപ്പോൾ}} \emph{then}.) Answer \emph{\ml{നിങ്ങൾ} \ml{എപ്പോൾ} \ml{പോകും}?} with this moment. (\textbf{\ml{ഇപ്പോൾ} \ml{ഞാൻ} \ml{പോകുന്നു}}.)
\end{tcolorbox}
\section[ippōḷ enna ōrmma]{(I go, I am going, right now) — from cold}

\label{lesson:ML-R104-present-recall}



\begin{quote}
Close the lessons before this one. Say \emph{I go}, then \emph{I am going}, then \emph{right now I am going}.
\end{quote}

\begin{grammarlens}[title={Grammar Lens: one form, and a word in front of it}]
\noindent
\begin{tabularx}{\linewidth}{@{}>{\raggedright\arraybackslash}X>{\raggedright\arraybackslash}X>{\raggedright\arraybackslash}X@{}}
\toprule
\textbf{} & \textbf{} & \textbf{} \\
\midrule
\textbf{\ml{ഞാൻ} \ml{പോകുന്നു}} & \emph{ñān pōkunnu} & I go \textbf{and} I am going \\
\textbf{\ml{ഇപ്പോൾ} \ml{ഞാൻ} \ml{പോകുന്നു}} & \emph{ippōḷ ñān pōkunnu} & right now I am going \\
\bottomrule
\end{tabularx}

\textbf{The verb is the same in both rows.} That is the whole chapter in one column. Where English changes \emph{go} to \emph{am going} and then adds \emph{right now}, Malayalam never touches the verb at all — it puts a word in the slot at the front, the same slot that has held \emph{today}, \emph{tomorrow}, a named day and a named hour.
\end{grammarlens}

\begin{grammarlens}[title={What you've built}]
\textbf{Not one word in this chapter was new.} The verb ending was named for you long ago; \textbf{\ml{ഇപ്പോൾ}} and \textbf{\ml{എപ്പോൾ}} have both been yours for chapters. What you did not have was permission to use one form for two English tenses, and the warning that goes with it.

\textbf{That warning is worth keeping.} English's \emph{am} has no home in \textbf{\ml{ഇപ്പോൾ} \ml{ഞാൻ} \ml{വരുന്നു}}. \textbf{\ml{ആണ്}} exists, and it joins a thing to what it \emph{is}. Propping up another verb is not its job — and the \emph{am going} is already inside \textbf{\ml{വരുന്നു}}, so there is nothing left for a second word to carry.
\end{grammarlens}

\subsection*{Your turn}
Say these aloud:

\begin{itemize}
\raggedright
  \item \emph{ñān pōkunnu}, meaning each of the two English sentences in turn
  \item \emph{ippōḷ ñān varunnu}
  \item \emph{eppōḷ}, then answer it
  \item where \emph{am} goes in Malayalam
\end{itemize}

\begin{tcolorbox}[breakable,colback=teal!4,colframe=teal!35!black,title={Before you move on}]
Two English sentences, one Malayalam one — which? (\emph{I go} and \emph{I am going}, both \textbf{\ml{ഞാൻ} \ml{പോകുന്നു}}.) How do you pin it to this moment? (\textbf{\ml{ഇപ്പോൾ}}, at the front.) Where does \emph{am} go? (\textbf{Nowhere} — it is inside the verb.)
\end{tcolorbox}
